% ==============================================================
% 第2章 关键技术介绍
% ==============================================================
\chapter{关键技术介绍}

"社区老年糖尿病多智能体随访管理系统"涉及多项关键技术，包括B/S系统架构、Django Web框架、大语言模型（LLM）、多智能体系统（MAS）与LangGraph框架、检索增强生成（RAG）技术、语音识别与语音合成技术等。本章将对上述技术进行介绍。

% ------------------------------------------------------------------
\section{系统架构——B/S}

B/S架构（Browser/Server架构）是主流的软件系统架构之一。与C/S架构不同，B/S架构将主要业务逻辑部署于服务器端，用户无需安装客户端软件，通过浏览器即可直接访问系统\todocite{}。这种架构的核心优势在于跨平台性和易维护性，用户通过浏览器实现多端使用，开发者只需更新服务器端的代码即可完成系统升级。

在医疗信息化领域，B/S架构已被广泛采用。医院的HIS（医院信息系统）、远程会诊系统以及社区卫生服务信息平台均普遍采用B/S架构，便于医务人员通过内网浏览器直接访问患者电子健康档案和管理系统\todocite{}。在慢性病管理场景中，B/S架构使老年患者无需安装专用APP，仅通过浏览器即可完成健康数据录入和查看健康报告，有效降低了技术使用门槛。

随着HTML5、WebGL等前端技术的发展以及云计算服务的普及，B/S架构的应用范围不断扩大，复杂的交互功能（如语音录制、音频播放、数据可视化等）均可在浏览器中流畅运行，为B/S架构提供了强大的技术支撑\todocite{}。本系统选用B/S架构，患者与医生均通过浏览器访问系统，保证了系统的可用性和部署便捷性。

% ------------------------------------------------------------------
\section{Web框架——Django}

Django是基于Python语言的高级Web框架，采用MVT（Model-View-Template）设计模式\todocite{}。Django以"全栈式"框架著称，内置了对象关系映射（ORM）、表单处理、模板引擎、Admin管理后台等组件，遵循"约定优于配置"的原则，能够帮助开发者快速构建功能完备的Web应用。该框架特别强调安全性，默认提供CSRF防护、SQL注入防护、XSS防护等安全机制。

在医疗信息系统开发中，Django展现出独特优势。其强大的ORM层能够屏蔽底层数据库差异，简化患者健康档案、随访记录、血糖数据等结构化信息的持久化操作；Admin后台为系统管理员提供开箱即用的数据管理界面，降低运维成本\todocite{}。此外，Django的用户认证系统（AbstractUser）支持灵活的角色扩展，便于实现患者与医生的角色隔离和权限控制。

随着Python在人工智能领域的广泛应用，Django常作为AI服务的Web接口框架，将LLM调用、多Agent编排等AI能力封装为标准的HTTP API，实现前后端的松耦合集成\todocite{}。本系统基于Django 5.2.13框架构建，利用MVT架构实现前后端分离，通过ORM管理SQLite/MySQL数据库的读写操作。

% ------------------------------------------------------------------
\section{大语言模型技术}

大语言模型（Large Language Model，LLM）是基于Transformer架构、在海量文本语料上预训练的生成式AI模型，具备文本生成、知识问答、逻辑推理、代码生成等多项能力\todocite{}。LLM通过自注意力（Self-Attention）机制捕捉文本序列中的长距离依赖关系，在上下文理解和复杂推理任务上展现出强大能力。

在医疗领域，LLM已被验证可应用于临床问答、病历摘要生成、药物相互作用分析、患者教育等多个场景。研究表明，GPT-4在美国医师执照考试（USMLE）中的表现已达到人类医学生的水平\todocite{}。然而，LLM在医疗场景中面临幻觉（Hallucination）问题——模型可能生成听起来合理但实际不准确的医疗建议，这要求在系统设计中引入知识库约束和人工审核机制。

本系统最终选用阿里云通义千问Qwen3.5系列作为核心LLM服务，通过阿里云DashScope平台的OpenAI兼容协议进行调用。选择Qwen3.5的核心原因在于：其一，中文理解与生成能力在国产大模型中处于领先水平，尤其在医疗术语和口语化表达的处理上表现优异；其二，DashScope平台同时提供LLM、ASR和TTS三类服务，使用统一的API Key即可调用，降低了系统集成复杂度；其三，数据存储于国内服务器，满足医疗数据合规性要求。系统同时通过统一的API封装层保留了切换至GPT-4o-mini（OpenAI）或GLM-4（智谱AI）的能力。系统通过角色设定（Role-Playing）、结构化输出约束（JSON Schema）和提示词工程（Prompt Engineering）等手段控制LLM在医疗场景中的行为准确性\todocite{}。

% ------------------------------------------------------------------
\section{多智能体系统与LangGraph框架}

\subsection{多智能体系统}

多智能体系统（Multi-Agent System，MAS）是由多个智能体（Agent）组成的协作计算系统，各Agent通过感知环境、自主推理和协调行动完成复杂任务\todocite{}。在医疗领域，MAS的核心优势在于通过专业化分工实现多维度决策覆盖：例如，MDAgents框架根据医疗任务复杂度动态分配单Agent独立决策或多Agent团队协作模式，在10个医疗基准测试中的7个上超越了现有方法\todocite{}；RareAgents框架通过多学科团队Agent的多轮讨论机制达成罕见病诊断共识，在诊断任务上优于GPT-4o。

本系统设计了五个专业化Agent：PatientAgent（负责健康数据解析与健康反馈生成）、TriageAgent（负责风险评估与预警分级）、SchedulerAgent（负责随访排程建议）、DoctorAgent（负责诊疗摘要生成）和MedicationAgent（负责用药方案审查），各Agent通过结构化状态传递实现协同决策，模拟了"全科医生——专科顾问——患者管理"的分层临床决策结构\todocite{}。

\subsection{LangGraph框架}

LangGraph是当前构建有状态多Agent应用的主流工程框架，基于LangChain扩展，采用有向图（Graph）的方式建模Agent工作流\todocite{}。其基本构成要素包括：图（Graph）定义工作流拓扑结构；节点（Node）代表单个Agent或功能模块；边（Edge）定义节点间信息传递方向；条件边（Conditional Edge）基于中间状态实现动态路由分支；状态（State）作为贯穿工作流的共享信息容器，确保各Agent能够访问完整上下文。

LangGraph与其他多Agent框架（如AutoGen、CrewAI）的核心区别在于其支持循环（Cycles）和递归结构的有状态工作流，天然适合慢病管理场景中需要迭代决策和长期记忆的应用需求\todocite{}。本系统基于LangGraph 1.1.6构建多Agent编排图，通过条件边实现"如果风险等级为红码则路由至紧急预警"等动态决策逻辑，状态机制保障了患者全链路上下文的一致性。

% ------------------------------------------------------------------
\section{检索增强生成技术}

检索增强生成（Retrieval-Augmented Generation，RAG）是将大语言模型与外部动态知识库相结合的技术范式，通过在生成阶段注入检索到的相关知识片段，有效解决LLM知识截止日期限制和领域知识深度不足的问题\todocite{}。

RAG的核心工作流程包括三个阶段：（1）文档预处理与向量化——将领域文档拆分为语义完整的文本块，通过嵌入模型（如all-MiniLM-L6-v2）将文本块转化为高维向量并存入向量数据库；（2）语义检索——当接收到用户查询时，将查询文本同样向量化，通过余弦相似度计算检索Top-K最相关的知识片段；（3）增强生成——将检索到的知识片段拼接至LLM提示词中，引导模型基于权威知识生成有依据的回答\todocite{}。

在糖尿病管理领域，RAG技术的引入具有特殊价值。一项实证研究表明，引入RAG框架后，GPT-4、Claude 2等模型在糖尿病教育场景中的准确回答率平均提升12\%\todocite{}。本系统以《中国糖尿病防治指南（2024版）》和《初级医疗2型糖尿病指南（2025版）》等权威文献为知识来源，采用ChromaDB 1.5.7向量数据库存储知识嵌入向量，通过all-MiniLM-L6-v2嵌入模型实现语义检索，为PatientAgent生成健康反馈建议时提供循证医学知识支撑，从根本上降低LLM的幻觉风险。

% ------------------------------------------------------------------
\section{语音识别与语音合成技术}

\subsection{ASR语音识别技术}

自动语音识别（Automatic Speech Recognition，ASR）技术负责将患者的口语化音频输入转化为可供系统处理的文本信息，是本系统适老化语音交互链路的首要环节。

现代ASR技术经历了从传统声学模型（HMM-DNN）到端到端序列模型（CTC、Attention-based Encoder-Decoder）再到大规模预训练模型的三次范式演进。OpenAI于2022年发布的Whisper模型是当前端到端ASR的代表性工作，其采用基于Transformer的Encoder-Decoder架构，在68万小时多语言弱监督数据上进行训练，在中文语音识别任务上展现出接近人类水平的准确率\todocite{}。阿里云达摩院推出的Paraformer系列模型则针对中文口语场景进行了深度优化，在医疗领域术语（如药品名称、检查指标等）的识别精度上相较通用ASR模型有显著提升\todocite{}。字节跳动火山引擎的流式语音识别服务同样支持中文实时转写，在低延迟与高并发场景下具有工程优势\todocite{}。

本系统的语音识别采用服务端ASR架构：前端通过浏览器MediaRecorder API录制音频，将音频文件上传至后端，后端调用ASR服务完成转写。系统以策略模式封装ASR服务调用接口，通过配置文件切换底层ASR供应商（Whisper API / 阿里云Paraformer / 火山引擎），实现ASR服务的供应商无关性，既降低了对单一供应商的技术依赖，也为后续根据实际部署环境选择成本最优方案提供了灵活性。

\subsection{TTS语音合成技术}

文本转语音（Text-to-Speech，TTS）技术负责将系统生成的文本信息（如健康反馈建议）转化为语音音频，以语音播报的形式呈现给患者，是本系统适老化无障碍设计的重要组成部分。

当代神经网络TTS系统通常由文本前端（Text Frontend）、声学模型（Acoustic Model）和声码器（Vocoder）三部分组成。微软Azure Neural TTS基于FastSpeech 2架构与HiFi-GAN声码器，提供超过400种高自然度语音音色，其中包括多种中文女声与男声选项。微软Edge TTS作为Azure Neural TTS的轻量级在线接口，无需注册即可使用，语音质量与Azure TTS一致，在本系统的适老化场景中可直接选用温和、语速较慢的中文女声音色（如XiaoxiaoNeural），有效降低老年患者的听觉理解负担\todocite{}。阿里云CosyVoice与火山引擎语音合成则分别提供具有情感表达能力的中文音色，支持语速、音调等参数调节\todocite{}。

本系统采用阿里云DashScope平台的CosyVoice语音合成服务作为默认TTS方案，选用温和女声音色（longxiaochun），语速适中、吐字清晰，符合老年患者的听觉习惯。同时通过策略模式预留了Edge TTS（微软免费方案）与火山引擎TTS的切换接口。合成结果以MP3格式缓存于服务端（基于voice+text的MD5哈希去重），前端通过HTML5 \texttt{<audio>}标签实现播放，避免重复合成相同内容带来的延迟与资源浪费。

% ------------------------------------------------------------------
\section{前端技术栈}

\subsection{HTML、CSS与JavaScript}

HTML、CSS和JavaScript是现代Web前端开发的核心技术组合。HTML（超文本标记语言）负责定义网页的结构和内容；CSS（层叠样式表）控制HTML元素的呈现样式，包括布局、颜色、字体等视觉效果；JavaScript作为脚本语言实现用户交互、动态内容更新等行为逻辑。三者各司其职又相互配合，是所有Web应用的技术基石\todocite{}。

\subsection{Bootstrap框架}

Bootstrap是基于HTML、CSS和JavaScript的前端框架，提供了一套预定义的样式组件和响应式栅格布局系统。其内置的按钮、导航栏、表单、模态框等UI组件通过CSS类实现，同时依赖JavaScript提供交互功能\todocite{}。Bootstrap的响应式设计保证系统在PC端、平板和手机端的一致呈现，其字体放大、高对比度配色等无障碍设计特性有助于满足老年用户的交互需求。本系统基于Bootstrap 5构建适老化前端界面，遵循《互联网应用适老化及无障碍改造指南》技术要求，采用暖色调配色方案（主色暖橙\#D35400，背景暖白\#FFF8F0），正文字号不小于18px，主要操作组件焦点区域不小于60$\times$60px，颜色对比度不低于4.5:1，核心操作流程控制在3步以内完成\todocite{}。

\subsection{Chart.js数据可视化}

Chart.js是一个简洁、灵活的JavaScript图表库，支持折线图、柱状图、饼图等常见图表类型\todocite{}。相比ECharts等重量级可视化工具，Chart.js以轻量级和响应式特性见长，适合需要在有限页面空间内展示健康趋势数据的场景。本系统在患者仪表盘中使用Chart.js展示血糖波动趋势折线图和风险等级分布饼图，帮助患者和医生直观了解健康数据的变化规律。
